\begin{abstract}
Let $G$ be the graph on the positive integers in which $u<v$ are adjacent exactly when $v/u$, written in lowest terms, has the form $(q+1)^2/q$ with $q\ge 1$. We study the symbolic dynamics of the \emph{paths of $G$ itself}: the nonbacktracking walks, regarded as a subshift over directed edges, with no rule selecting among the admissible continuations. The graph is triangle-free with girth $4$ and odd girth $5$; every closed walk yields a multiplicative relation among the numbers $(N+1)^2/N$, and cycles correspond to those relations that are realised by admissible words with distinct partial products. For every integer $m\ge1$ the graph contains an explicit finite subgraph $F_m$ consisting of a hub edge $x\!-\!4x$, $x=2^m$, and $m$ internally disjoint paths of length $3$ between its endpoints. We prove that the nonbacktracking entropy of $F_m$ is exactly $\log\lambda_m$, where $\lambda_m$ is the real root of $\lambda^3-\lambda^2+\lambda=m$, so that $\tfrac13\log m<\hNB(F_m)<\tfrac13\log m+\tfrac{1}{3}m^{-1/3}$ for $m\ge2$; we give the complete spectrum of the edge matrix, the Ihara-type determinant, the numbers of periodic symbolic paths of every period, and the maximal-entropy Markov measure with its suspension coding. Since $F_m$ is also the induced subgraph of $G$ on its vertex set, the same value is the entropy of the induced window. Consequently the finite-subsystem capacity of $G$ is infinite, and $\hNB(G[1,X])\to\infty$. Strictly increasing (forward) walks have bounded length in every finite subgraph, and any deterministic finite-state selection of paths has zero entropy: the entropy of $G$ is a property of reversible histories, and it coexists with the zero-entropy Sturmian behaviour of the prime-valuation selector and the aperiodic forward arithmetic.
\end{abstract}

\tableofcontents
\bigskip

\section{Introduction}\label{sec:intro}

The arithmetic graph $G$ studied here has vertex set $\N=\{1,2,3,\dots\}$ and edges
\[
\{kN,\;k(N+1)^2\},\qquad k,N\ge1 .
\]
Equivalently, two integers $u<v$ are adjacent iff $v/u=(q+1)^2/q$ in lowest terms for some $q\ge1$, in which case $q=N$ and $k=u/q$. Moving along the edge of label $N$ multiplies by $(N+1)^2/N$ going up and divides by it going down; the up-move $U_N$ requires $N\mid x$ and the down-move $D_N$ requires $(N+1)^2\mid x$. The graph is connected, locally finite, and has a rich cycle structure; its vertex $x$ has degree $\tau(x)+\#\{d\ge2:d^2\mid x\}$ (Proposition~\ref{prop:degree}; \cite{UTI}). Connectivity follows from the bounded successor distance of~\cite{UBND}.

\paragraph{Results.}
\begin{itemize}[leftmargin=1.6em,itemsep=2pt]
\item \textbf{Cycles} (Theorem~\ref{thm:cycles}, Proposition~\ref{prop:relations}; triangle-freeness and odd girth $5$: \cite{UTHREE}; girth $4$: the correction \cite{UCOR}). $G$ is triangle-free, has girth~$4$ and odd girth~$5$. A closed walk with step labels $(N_i,\varepsilon_i)$ forces the relation $\prod_i r_{N_i}^{\varepsilon_i}=1$ for $r_N=(N+1)^2/N$; conversely such a relation gives a closed walk exactly when the word is admissible at the base point. Triangle-freeness says that $r_ar_b=r_c$ has no solution in positive integers, and the shortest odd cycle realises the relation $r_1^2r_6=r_2r_{27}$.
\item \textbf{Exact entropy} (Theorems~\ref{thm:entropy} and~\ref{thm:induced}; proved in this paper, the inherited release~\cite{UENT} giving only the lower bound $\tfrac14\log(\tau(x)-1)$). Let $F_m\subseteq G$ be the subgraph with vertices $x,4x,y_j,4y_j$ ($x=2^m$, $y_j=2^{m-j}(2^j+1)^2$, $1\le j\le m$) and edges $x\!-\!4x$ and $x\!-\!y_j\!-\!4y_j\!-\!4x$. Then
\[
\begin{gathered}
\hNB(F_m)=\log\lambda_m,\qquad \lambda_m^3-\lambda_m^2+\lambda_m=m,\\
\tfrac13\log m<\log\lambda_m<\tfrac13\log m+\tfrac{1}{3m^{1/3}}\quad(m\ge2),
\end{gathered}
\]
and the induced subgraph of $G$ on $V(F_m)$ is $F_m$ itself.
\item \textbf{Spectrum and periodic paths} (Theorems~\ref{thm:spectrum} and~\ref{thm:periodic}; proved in this paper). The edge matrix has characteristic polynomial $(\lambda^6-1)^{m-1}(\lambda^2-1)(\lambda^3-\lambda^2+\lambda-m)(\lambda^3+\lambda^2+\lambda+m)$; $F_m$ has exactly $m$ four-cycles and $\binom m2$ six-cycles, and the number of closed nonbacktracking walks of length $n$ is $6(m-1)[6\mid n]+(1+(-1)^n)(1+s_n)$ with $s_n$ the power sums of the roots of $\lambda^3-\lambda^2+\lambda-m$.
\item \textbf{Capacity} (Theorem~\ref{thm:capacity}; the unboundedness of finite-subsystem entropy is first obtained, with the weaker bound $\tfrac14\log b(x)$, in~\cite{UENT}; the exact value is proved here). For every $x$, $F(x)\subseteq G$ is a theta graph with $\hNB(F(x))=\log\lambda_{\tau(x)-1}$. Hence the finite-subsystem capacity $\Hfin(G)=+\infty$, and $\hNB(G[1,X])\ge\tfrac13\log(\tau(x)-1)$ whenever $4(x+1)^2\le X$.
\item \textbf{Three dynamics} (Section~\ref{sec:three}). Forward arithmetic, autonomous prime-valuation selection, and reversible graph histories are three different dynamical systems on the same arithmetic; their entropies ($0$ for the first two in the sense of Proposition~\ref{prop:forward}, unboundedly large for the third) do not conflict.
\end{itemize}

\paragraph{Provenance.} The graph and its triangle-freeness and odd girth $5$ come from the frozen release \texttt{UNNS\_U3}~\cite{UTHREE}; the value $\operatorname{girth}(G)=4$ is the correction of \texttt{UNNS\_U3\_COR}~\cite{UCOR}; the degree formula is in \texttt{UNNS\_U2\_TI\_DYN}~\cite{UTI}; connectivity follows from the successor bound of \texttt{UNNS\_U1\_BND}~\cite{UBND}. The memory-two metric shift $Y_{\mathrm{NB}}$ is defined in \texttt{UNNS\_U4\_STAGE2}~\cite{USTAGE}; the family $F(x)$ and the bound $\hNB(F(x))\ge\tfrac14\log(\tau(x)-1)$ come from \texttt{UNNS\_U4\_M2\_ENT}~\cite{UENT}; the Sturmian selector $\MAXPP$ and the valuation law come from \texttt{UNNS\_PRIME\_DYN} (current Version 1.2; the archive held here is Version 1.1)~\cite{PD}. The present paper gives self-contained proofs of everything it uses from the first group, replaces the inherited lower bound by the exact value, and reconciles the new formula with the inherited numerical entropies (Section~\ref{sec:recon}). The exact entropy formula, the spectrum, the Ihara-type determinant, the periodic-path counts, the maximal-entropy measure, the spectral gap and the induced-subgraph theorem are results of this paper. References are listed with their SHA-256 hashes at the end, and Appendix~\ref{app:prov} gives the source of each theorem family; all statements about the new mathematics are verified by the independent script described in Appendix~\ref{app:ver}.

\paragraph{Organisation.} Section~\ref{sec:graph} sets up the path grammar of $G$. Section~\ref{sec:nb} fixes conventions for nonbacktracking histories. Section~\ref{sec:cycles} treats cycles and their relation to multiplicative identities. Section~\ref{sec:theta} proves the exact entropy and the spectrum of theta graphs, with comparison results and a complete worked example. Section~\ref{sec:F} applies this to $F(x)$ and $F_m$, proves the induced-subgraph theorem and works out $F_2$, $F_3$ and $F(12)$. Section~\ref{sec:symbolic} gives the transition matrices, periodic orbits, the maximal-entropy measure with its stationary distribution, and the spectral gap. Section~\ref{sec:capacity} proves the capacity theorem and quantifies it on finite windows. Section~\ref{sec:three} compares the three dynamics, Section~\ref{sec:recon} reconciles with U4-M2-ENT, and Section~\ref{sec:synthesis} collects the consequences.
